\documentclass[11pt]{article}
\usepackage{amsmath,amssymb,amsthm}
\newtheorem{theorem}{Theorem}
\title{A Quadratic Pole Budget for Compact Disjoint Dominance}
\author{}
\date{November 24, 2026}
\begin{document}
\maketitle

\paragraph{Status and scope.}
This is a complete written argument, with a small exact companion check,
not a Lean theorem. It concerns a particular sign-cube mechanism: each
pole belongs to one independently signed group, and the output word is
that group-sign word up to a fixed signed permutation. It does not bound
general nonlinear sign-cube encodings, adaptive residues, the full local
compiler, or worst-case head complexity. No model definition is changed.

Let $q\geq1$. For $j=1,\ldots,q$, let
\[
 f_j(t)=\sum_{h=1}^{m_j}\frac{c_{jh}}{t-p_{jh}},
 \qquad m_j\geq1,\qquad G=\sum_jm_j.
\]
Repeated poles and arbitrary real residues are allowed. Suppose all poles
belong to $[0,a]$, and $t_1,\ldots,t_q\geq b>a$.
The independent group signs $x_j\in\{-1,1\}$ produce scores
$\sum_jx_jf_j(t_i)$.

After relabeling and changing fixed group signs, the claim that the
strict output word is exactly $x$ for every $x$ is equivalent to
\[
 f_i(t_i)>\sum_{j\ne i}|f_j(t_i)| \qquad(1\leq i\leq q).
 \tag{1}
\]
Indeed, the minimum of $x_i\sum_jx_jf_j(t_i)$ over all sign choices is
$f_i(t_i)-\sum_{j\ne i}|f_j(t_i)|$. Thus this assumption is both
necessary and sufficient for the stated fixed signed-permutation
decoder. It is not necessary for a general bijective threshold map.

\begin{theorem}[Finite compactness inequality]
Put $L=\lfloor 2G/q\rfloor$. Under (1),
\[
 \frac q2<L\left(1+(1-a/b)^{-2L}\right).
 \tag{2}
\]
More precisely, the number $s$ of groups with $m_j\leq L$ satisfies
$s>q/2$ and
\[
 s<L\left(1+(1-a/b)^{-2L}\right).
 \tag{3}
\]
\end{theorem}
\begin{proof}
Fewer than $q/2$ groups can have $m_j>2G/q$, since their total size
would otherwise exceed $G$. Let $S$ be the remaining groups, and restrict
both rows and columns of (1) to $S$. Strict diagonal dominance persists.

For $j\in S$ put
\[
 Q_j(t)=t^{L-m_j}\prod_{h=1}^{m_j}(t-p_{jh}),
 \qquad P_j(t)=Q_j(t)f_j(t).
\]
Every $P_j$ is a polynomial of degree at most $L-1$, including with
repeated poles. The factors $t^{L-m_j}$ are an algebraic padding only;
they introduce no heads and no assertion of an admissible zero pole.
For $t\geq b$,
\[
 t^L(1-a/b)^L\leq Q_j(t)\leq t^L.
 \tag{4}
\]
Consider the $s\times s$ real matrix
\[
 R_{ij}=\frac{P_j(t_i)}{P_i(t_i)}
 =\frac{f_j(t_i)}{f_i(t_i)}\frac{Q_j(t_i)}{Q_i(t_i)}.
\]
Its denominators are positive by (1) and (4). Its diagonal entries are
one, its rank is at most $L$, and every off-diagonal row absolute sum
is less than $M=(1-a/b)^{-L}$. Consequently
\[
 \|R\|_F^2
 =s+\sum_i\sum_{j\ne i}|R_{ij}|^2
 \leq s+\sum_i\left(\sum_{j\ne i}|R_{ij}|\right)^2
 <s(1+M^2).
\]
For any real square matrix of rank $k$,
$|\operatorname{tr}R|\leq\sqrt{k}\|R\|_F$.
For completeness, write a singular value decomposition
$R=U\operatorname{diag}(\sigma_1,\ldots,\sigma_k)V^\top$.
The trace has absolute value at most $\sum\sigma_i$, which is at most
$\sqrt{k}(\sum\sigma_i^2)^{1/2}$ by Cauchy--Schwarz.
Here $\operatorname{tr}R=s$, so $s^2<Ls(1+M^2)$, proving (3).
\end{proof}

\begin{theorem}[Quadratic pole cost for a compact germ]
Fix $q,m_1,\ldots,m_q$. Allow all residues, poles and query times to
depend on a positive parameter. Suppose (1) holds for arbitrarily small
parameter values, with $a/b\to0$. Then
\[
 G\geq q^2/8.
 \tag{5}
\]
\end{theorem}
\begin{proof}
The integer $L$ is fixed. Taking the limit in (2) gives
$q/2\leq2L\leq4G/q$.
\end{proof}

\paragraph{A finite quantitative version.}
If $a/b\leq1/(4L)$, Bernoulli's inequality gives
$(1-a/b)^L\geq3/4$. Equation (2) therefore implies
\[
 G>9q^2/100.
\]
For an order-$\nu$ kernel $(t-p)^{-\nu}$ the identical proof pads to
degree $\nu L$, and gives $G\geq q^2/(8\nu)$ in the compact limit.

\paragraph{Why the compactness assumption matters.}
Outside the compact regime, a linear pole budget can give independent
dominance. Fix $R>1$, $a_j=R^j$, $t_i=a_i$, and
\[
 f_j(t)=\frac{a_jt}{(t+a_j)(t+2a_j)}
       =\frac{-a_j}{t+a_j}+\frac{2a_j}{t+2a_j}.
\]
The diagonal value is $1/6$. For $i\ne j$,
$0<f_j(t_i)\leq\min\{a_i/a_j,a_j/a_i\}$, so every off-diagonal
row sum is at most $2/(R-1)$. Taking $R=32$ proves (1) with $2q$ poles.
These poles and query times occupy matching geometric scales. This is
not a compact reverse-cube construction: the required full-cube slope
budget must still be established after assigning tokens to slots.
Applying a different projective transform to each slot does not by
itself establish that budget when a slot receives many token scales
across different represented rows.

\paragraph{An exact projective-budget obstruction for those scales.}
Here is a precise limitation of that proposed transfer. Suppose one
slot receives two different poles $P_0<P_1$, and its query list contains
two values $T_0<T_1$, with $P_1<T_0$, each repeated at least $L$ times.
Consider a projective change of coordinate after which all query
coordinates lie strictly above all pole coordinates. Assume the
transformed differences are extended to a reverse affine denominator
on the full row/query Boolean cube. Then necessarily
\[
 \frac{(P_1-P_0)(T_1-T_0)}{(T_0-P_0)(T_1-P_1)}<\frac1L.
 \tag{6}
\]
For clarity, this is an obstruction to the stated headwise projective
transfer of this fixed denominator grid. It is not an obstruction to
an alternative representation of its signs.

To prove (6), work first in transformed coordinates, denoting the
two displayed poles by $x_0<x_1$ and the two queries by $y_0<y_1$.
An affine denominator extending the differences has row slopes
$\alpha_i=x_i-u$, column slopes $\beta_j=v-y_j$, and constant
$v-u$, for some $u,v$. Nonnegative slopes imply
$u\leq\min_i x_i$ and $v\geq\max_j y_j$. The full-cube budget
therefore implies
\[
 v-u>\sum_i\alpha_i+\sum_j\beta_j
 \geq (x_1-u)+L(v-y_0).
\]
Writing $\Delta P=x_1-x_0$, $\Delta T=y_1-y_0$ and $g=y_0-x_1$,
this gives $g>(L-1)(v-y_0)\geq(L-1)\Delta T$, and hence
\[
 \frac{\Delta P\Delta T}{(g+\Delta P)(g+\Delta T)}
 <\frac{\Delta T}{g+\Delta T}<\frac1L.
\]
Projective transformations preserve this cross-ratio. If the chart has its pole in
the separating gap, it reverses the orders within both intervals;
the same cross-ratio is obtained by reversing both pairs. The
multiplicity assumption on both displayed queries covers this case.
A chart cannot split either pole or query interval while preserving
strict separation of every transformed pole and query.

For the geometric pool, take poles $-a_k,-a_j$ and query values
$a_j,a_k$, where $k>j$ and $a_k/a_j\geq32$. The left side of (6) is
at least $(31/33)^2>1/2$. Hence a slot mixing even two such scales
cannot satisfy this reverse extension budget if each address query
is used twice. In a signed color-mask assembly a slot of mask weight
at least two has precisely that repetition. Thus arbitrary token
permutation across these geometric scales is not licensed by
per-slot projective compactification.

\paragraph{Relation to the four-query example.}
The saved eight-pole rectangular certificate has a sufficient sign
subcube
\[
 (\tau_0,\ldots,\tau_7)=(a,b,c,d,-c,c,-a,a).
\]
Its output word is $(d,-a,b,c)$. Thus its disjoint group sizes are
$3,1,3,1$, and it lies within the scope of this note after the stated
signed permutation. This observation does not imply that every
successful rectangular sign-cube family has such a subcube.

\paragraph{Verification boundary.}
The companion script checks that sixteen-word identity using exact
rationals when the saved pool is available, verifies the rational
padding identity on small examples, and checks the noncompact
two-pole-per-group construction exactly. Those tests audit identities;
the rank argument above supplies the unrestricted quantitative claim
within its explicit scope.
\end{document}
